\chapter{Zeta Functions and Densities}
\begin{center}\itshape
The field $\Q(\sqrt{-23})$ has class number $3$.\\
The series $1-\frac12-\frac13+\frac14-\frac15+\frac16-\frac17+\frac18+\frac19-\cdots$, signs given by the squares modulo $23$, sums to $\frac{3\pi}{\sqrt{23}}$.\\
Coincidence?
\end{center}

In 1839 Dirichlet published a long memoir with a modest title, \emph{Recherches sur diverses applications de l'analyse
infinit\'esimale \`a la th\'eorie des nombres}. Two years earlier he had used infinite series to prove that every arithmetic
progression contains primes (Volume~5). Now he turned the same weapon on a question of pure algebra: how many classes of
binary quadratic forms of discriminant $D$ are there? Gauss had computed long tables and found no law. Dirichlet found a
formula, and it was a shock: the class number, a finite count of algebraic objects, equals an infinite series times an
elementary constant. For $D=-23$ the series above, multiplied by $\frac{\sqrt{23}}\pi$, gives exactly $3$. A message from
algebra, written in the language of analysis.

Kummer understood the message at once. In 1847--1850, fighting with Fermat's equation (Chapter~5), he applied Dirichlet's
method to $\Q(\zeta_p)$ and found that the class number is controlled by Bernoulli numbers, those strange rationals that
Jakob Bernoulli had met while summing powers of integers. That is how $37$ was unmasked as the first irregular prime: $37$
divides the numerator of $B_{32}=-\frac{7\,709\,321\,041\,217}{510}$. In the supplements he wrote for Dirichlet's lectures, Dedekind gave every number field its zeta
function, and a few years later Kronecker and Frobenius used it to ask how primes split. Frobenius conjectured in 1880, and
dared to publish only in 1896, that in a Galois extension each conjugacy class $C$ of the Galois group is the Frobenius class
of a proportion $\frac{|C|}{|G|}$ of all primes. The proof came in 1922 from Nikolai Chebotarev, a young mathematician
working in Odessa in the hungry years after the civil war, with an idea so ingenious that Artin used it five years later to
prove his reciprocity law, the heart of class field theory (Chapter~10).

This chapter follows that road. The analytic engine is a single counting theorem: the number of ideals of norm at most $x$
in an ideal class is $\kappa x+O(x^{1-1/n})$, with an explicit constant $\kappa$ that contains the regulator, the
discriminant and $\pi$. It comes from counting lattice points in a cone, the geometry of Chapters~3 and~4. Everything else
follows: the analytic class number formula, the class numbers of cyclotomic fields and Kummer's criterion (closing the
promise of Chapter~5), and, with ray classes and Chebotarev's crossing trick, the density theorem announced in
Figure~\ref{fig:cycletypes}. We follow the proof of Lenstra and Stevenhagen, which needs no more analysis than Dirichlet's.

\begin{reminder}[from Volumes 2, 4, 5 and Chapters 2--5]
\emph{Dirichlet series} (Volume~4, Chapter~10): absolute convergence, Euler products (Volume~4, Theorem~10.3),
and $\zeta(s)-\frac1{s-1}$ holomorphic on $\operatorname{Re}s>0$ (Volume~4, Proposition~10.6).
\emph{Abel summation} (Volume~5, Lemma~1.1) and \emph{holomorphic integrals} (Volume~4, Theorem~9.15).
\emph{Characters of a finite abelian group $A$} (Volume~2, Section~11.4): there are $|A|$ homomorphisms
$\psi:A\to\C^\times$, and $\sum_\psi\psi(a)=|A|$ or $0$ according as $a=1$ or not. \emph{Dirichlet $L$-functions} (Volume~5,
Chapter~3): Gauss sums $|\tau(\chi)|=\sqrt q$ for primitive $\chi$ modulo $q$ (Volume~5, Theorem~3.5), and for odd
primitive $\chi$, $L(1,\chi)=-\frac{i\pi}{q\,\tau(\overline\chi)}\sum_{a=1}^{q-1}\overline\chi(a)a$ (Volume~5,
Theorem~3.10). \emph{Bernoulli numbers} (Volume~4, Proposition~1.14 and Theorem~1.15):
$B_n=0$ for odd $n\ge3$, and $\sum_{a=0}^{N-1}a^k=\frac1{k+1}\sum_{j=0}^k\binom{k+1}jB_jN^{k+1-j}$. From this book: unique
factorisation of ideals and the norm (Chapter~2), the Minkowski embedding $j$ and $\operatorname{covol}j(I)=2^{-r_2}
\sqrt{|d_K|}\,N(I)$ (Proposition~\ref{prop:covol}), the log embedding $\ell$ and the regulator (Chapter~4), Frobenius
elements (Theorem~\ref{thm:decomp}).
\end{reminder}

\section{The Dedekind zeta function}

Let $K$ be a number field of degree $n$. The \emph{Dedekind zeta function} of $K$ is
\[
\zeta_K(s)=\sum_{I}\frac1{N(I)^s},
\]
the sum over the non-zero ideals $I$ of $\mathcal O_K$. For $K=\Q$ it is Riemann's $\zeta$. Throughout, $\mathfrak p$ runs over the
non-zero prime ideals of $\mathcal O_K$, and $s=\sigma+it$.

\begin{proposition}\label{prop:zetaK}
For $\sigma>1$ the series converges absolutely, $\zeta_K(s)=\prod_{\mathfrak p}(1-N\mathfrak p^{-s})^{-1}$, $|\zeta_K(s)|\le\zeta(\sigma)^n$,
and
\[
\zeta_K(s)=\exp\Bigl(\sum_{\mathfrak p}\sum_{k\ge1}\frac{N\mathfrak p^{-ks}}k\Bigr)\ne0.
\]
\end{proposition}
\begin{proof}
Above each prime number $p$ lie at most $n$ primes $\mathfrak p$, each of norm $N\mathfrak p=p^{f}\ge p$ (Theorem~\ref{thm:sumef}).
For $y\ge2$, expanding the finite product of geometric series and using unique factorisation (Theorem~\ref{thm:ufi}) and the
multiplicativity of the norm (Proposition~\ref{prop:normideal}),
\[
\prod_{N\mathfrak p\le y}\bigl(1-N\mathfrak p^{-\sigma}\bigr)^{-1}=\sum_{I:\ \mathfrak p\mid I\Rightarrow N\mathfrak p\le y}N(I)^{-\sigma}
\]
contains every ideal of norm at most $y$, and is at most $\prod_{p\le y}(1-p^{-\sigma})^{-n}\le\zeta(\sigma)^n$. So the partial sums of
$\sum_IN(I)^{-\sigma}$ are bounded by $\zeta(\sigma)^n$, and the series converges absolutely. The Euler product and the exponential
formula follow exactly as for Dirichlet $L$-functions (Volume~5, Theorem~3.6(a)).
\end{proof}

The rest of the chapter is about the behaviour of $\zeta_K$ near $s=1$, and it all comes from counting ideals.

\section{Lattice points in cones}

How many points of a lattice $\Lambda\subseteq\R^n$ lie in a large dilate $tS$ of a set $S$? About $\operatorname{vol}(S)t^n/
\operatorname{covol}(\Lambda)$, if the boundary of $S$ is not too wild. Here is the version we need, for translated lattices.

\begin{lemma}[the Lipschitz principle]\label{lem:lipschitz}
Let $\Lambda\subseteq\R^n$ be a lattice and $S\subseteq\R^n$ a bounded measurable set whose boundary is contained in a finite
union of images $\phi_i([0,1]^{n-1})$ of Lipschitz maps $\phi_i:[0,1]^{n-1}\to\R^n$. Then for $t\ge1$ and every $v\in\R^n$,
\[
\#\bigl((v+\Lambda)\cap tS\bigr)=\frac{\operatorname{vol}(S)}{\operatorname{covol}(\Lambda)}t^n+O(t^{n-1}),
\]
with an implied constant depending on $\Lambda$, $S$ and the $\phi_i$, but not on $v$ or $t$.
\end{lemma}
\begin{proof}
Let $F$ be a fundamental parallelotope of $\Lambda$ and $\delta$ its diameter. The tiles $T_\lambda=v+\lambda+F$, $\lambda\in\Lambda$,
partition $\R^n$, and $T_\lambda$ contains the point $v+\lambda$. Let $a$ be the number of tiles contained in $tS$ and $b$ the
number of tiles meeting $tS$. Then $a\operatorname{covol}(\Lambda)\le t^n\operatorname{vol}(S)\le b\operatorname{covol}(\Lambda)$,
and the number of points of $v+\Lambda$ in $tS$ lies between $a$ and $b$. A tile that meets $tS$ without being contained in it
meets the boundary $t\,\partial S$, because a tile is connected. So it suffices to show that $O(t^{n-1})$ tiles meet $t\,\partial S$.

Let $c$ be a Lipschitz constant for all the $\phi_i$, and $k=\lceil t\rceil$. Cut $[0,1]^{n-1}$ into $k^{n-1}$ cubes of side $\frac1k$.
The image under $t\phi_i$ of each small cube has diameter at most $tc\sqrt{n-1}/k\le c\sqrt{n-1}=:d$, so it lies in a ball of
radius $d$, and every tile meeting it lies in the concentric ball of radius $d+\delta$. The tiles are disjoint of volume
$\operatorname{covol}(\Lambda)$, so there are at most $\operatorname{vol}(B_{d+\delta})/\operatorname{covol}(\Lambda)$ of them. In
total at most a constant times $k^{n-1}\le(2t)^{n-1}$ tiles meet $t\,\partial S$.
\end{proof}

Now the cones. Let $K$ have signature $(r_1,r_2)$, $r=r_1+r_2-1$, and use the notation of Chapter~4: $e_i=1$ for the real
coordinates and $e_i=2$ for the complex ones. On $K_\R=\R^{r_1}\times\C^{r_2}$, let $K_\R^\times$ be the set of points with all
coordinates non-zero, and for $x\in K_\R^\times$ put
\[
N(x)=\prod_{i=1}^{r+1}|x_i|^{e_i},\qquad \ell(x)=\bigl(e_1\log|x_1|,\dots,e_{r+1}\log|x_{r+1}|\bigr)\in\R^{r+1},
\]
so that $N(j(\alpha))=|N(\alpha)|$ and $\ell(j(u))=\ell(u)$ for $\alpha\in K^\times$, $u\in\mathcal O_K^\times$, and the coordinates of
$\ell(x)$ add up to $\log N(x)$. Let $\mathbf e=(e_1,\dots,e_{r+1})$, a vector whose coordinates add up to $n$.

Let $U\subseteq\mathcal O_K^\times$ be a subgroup of finite index. By the unit theorem (Theorem~\ref{thm:dirichlet}), $U/(U\cap\mu_K)$
is free of rank $r$; choose $\varepsilon_1,\dots,\varepsilon_r\in U$ whose classes form a basis, let $R_U$ be the absolute value of
an $r\times r$ minor of the matrix of the $\ell(\varepsilon_k)$ (the \emph{regulator of $U$}; it is $[\mathcal O_K^\times:U\mu_K]R_K$
by the argument of Proposition~\ref{prop:regwelldef}), and $w_U=|U\cap\mu_K|$. Let $P_U=\{\sum_kt_k\ell(\varepsilon_k):t\in[0,1)^r\}$,
the fundamental parallelotope of $\ell(U)$, and define
\[
D_U=\bigl\{x\in K_\R^\times:\ \ell(x)\in\R\,\mathbf e+P_U\bigr\},\qquad D_U(1)=\{x\in D_U:N(x)\le1\}.
\]

\begin{lemma}\label{lem:cone}
\begin{enumerate}[label=(\alph*)]
\item $D_U$ is a cone ($\lambda D_U=D_U$ for $\lambda>0$), stable under changing the signs of real coordinates and multiplying complex
coordinates by numbers of modulus~$1$.
\item For every $x\in K_\R^\times$, exactly $w_U$ elements $u\in U$ satisfy $j(u)x\in D_U$.
\item $D_U(1)$ is bounded, its boundary is contained in finitely many Lipschitz images of $[0,1]^{n-1}$, and
$\operatorname{vol}D_U(1)=2^{r_1}\pi^{r_2}R_U$. For a fixed choice of signs $s\in\{\pm1\}^{r_1}$, the part $D_U(1)_s$ where the real
coordinates have signs $s$ has the same properties, with volume $\pi^{r_2}R_U$.
\end{enumerate}
\end{lemma}
\begin{proof}
(a) $\ell(\lambda x)=\ell(x)+(\log\lambda)\mathbf e$, and $\ell$ only sees absolute values.

(b) The vectors $\mathbf e$ and $\ell(\varepsilon_1),\dots,\ell(\varepsilon_r)$ form a basis of $\R^{r+1}$: the $\ell(\varepsilon_k)$ form a basis of
the hyperplane $H$ of vectors with coordinate sum $0$ (unit theorem), and $\mathbf e\notin H$. Write $\ell(x)=c\,\mathbf e+\sum
y_k\ell(\varepsilon_k)$. Then $\ell(j(u)x)=\ell(u)+\ell(x)$, and $j(u)x\in D_U$ if and only if $\ell(u)=\sum m_k\ell(\varepsilon_k)$ with
$m_k=-\lfloor y_k\rfloor$. The units of $U$ with this value of $\ell$ form a coset of $\ker\ell\cap U=U\cap\mu_K$
(Proposition~\ref{prop:logkernel}), which has $w_U$ elements.

(c) Fix signs $s$ and define $\Phi:[0,1]\times[0,1]^r\times[0,1]^{r_2}\to K_\R$ by
\[
\Phi(\rho,t,\theta)_i=s_i\,\rho\,e^{\sum_kt_k\ell_i(\varepsilon_k)}\ (i\le r_1),\qquad
\Phi(\rho,t,\theta)_{r_1+j}=\rho\,e^{\frac12\sum_kt_k\ell_{r_1+j}(\varepsilon_k)}e^{2\pi i\theta_j}\ (j\le r_2).
\]
For $\rho>0$, $\ell(\Phi(\rho,t,\theta))=(\log\rho)\mathbf e+\sum t_k\ell(\varepsilon_k)$ and $N(\Phi)=\rho^n$, so $\Phi$ maps
$\left]0,1\right]\times[0,1)^r\times[0,1)^{r_2}$ bijectively onto $D_U(1)_s$. The map $\Phi$ is $C^1$ on a neighbourhood of the compact box,
hence bounded and Lipschitz on it (mean value inequality). On the open box it is injective, and its inverse $x\mapsto
(N(x)^{1/n},t(x),\theta(x))$, given by the coordinates of $\ell(x)$ in the basis above and by the arguments of the complex
coordinates, is continuous on an open set of $K_\R$ that is the image of the open box; so that image is open. Hence the boundary of
$D_U(1)_s$ is contained in $\Phi(\text{box})\setminus\Phi(\text{open box})\subseteq\Phi(\partial\,\text{box})$, and the boundary of the box is
made of $2(1+r+r_2)=2n$ faces, each a copy of $[0,1]^{n-1}$.

For the volume, use polar coordinates in each complex coordinate (Volume~6, Chapter~8) and then the variables $u_i=e_i\log|x_i|$:
$\dd x_i=e^{u_i}\dd u_i$ for a real coordinate of fixed sign, and $|z|\dd|z|\dd\varphi=\frac12e^{u}\dd u\dd\varphi$ for a complex one
with $u=2\log|z|$. So
\[
\operatorname{vol}D_U(1)_s=(2\pi)^{r_2}2^{-r_2}\int_\Omega e^{u_1+\dots+u_{r+1}}\dd u,
\]
where $\Omega=\{c\,\mathbf e+v:c\le0,\ v\in P_U\}$.
Change variables to $(c,t)$. The Jacobian is the absolute value of
\[
\det\bigl(\mathbf e,\ell(\varepsilon_1),\dots,\ell(\varepsilon_r)\bigr);
\]
adding all rows to the first one turns it into $(n,0,\dots,0)$, so the determinant is $\pm n$ times an $r\times r$ minor, that is
$nR_U$ in absolute value. As $u_1+\dots+u_{r+1}=nc$,
\[
\operatorname{vol}D_U(1)_s=\pi^{r_2}\,nR_U\int_{-\infty}^0e^{nc}\dd c=\pi^{r_2}R_U.
\]
The $2^{r_1}$ sign patterns give $2^{r_1}\pi^{r_2}R_U$.
\end{proof}

\begin{theorem}[counting in a cone]\label{thm:conecount}
Let $U\subseteq\mathcal O_K^\times$ have finite index, $\mathfrak a$ a non-zero fractional ideal, $\beta\in K$, and let $\Gamma$ be the
set of non-zero $\gamma\in\beta+\mathfrak a$, either all of them or only those whose real conjugates have prescribed signs $s$.
Suppose $U\Gamma=\Gamma$. Then the number of $U$-orbits in $\{\gamma\in\Gamma:|N(\gamma)|\le X\}$ is
\[
\frac{V}{w_U\operatorname{covol}j(\mathfrak a)}\,X+O\bigl(X^{1-1/n}\bigr),\qquad V=2^{r_1}\pi^{r_2}R_U\ \text{or}\ \pi^{r_2}R_U
\]
in the two cases.
\end{theorem}
\begin{proof}
Multiplication by $U$ preserves $|N|$ and $\Gamma$, and by Lemma~\ref{lem:cone}(b) each orbit has exactly $w_U$ points in
$j^{-1}(D_U)$. By the cone property, $j(\gamma)\in D_U$ and $|N(\gamma)|\le X$ if and only if $j(\gamma)\in X^{1/n}S$, with
$S=D_U(1)$ or $D_U(1)_s$. So the count is $\frac1{w_U}\#\bigl((j(\beta)+j(\mathfrak a))\cap X^{1/n}S\bigr)$, up to the possible point
$0$, and Lemma~\ref{lem:lipschitz} with $t=X^{1/n}$ and Lemma~\ref{lem:cone}(c) finish the proof.
\end{proof}

\section{Counting ideals}

\begin{theorem}[ideals in a class]\label{thm:idealcount}
Let $C$ be an ideal class of $K$. Then
\[
\#\{I\in C:\ I\subseteq\mathcal O_K,\ N(I)\le x\}=\kappa_Kx+O\bigl(x^{1-1/n}\bigr),\qquad
\kappa_K=\frac{2^{r_1}(2\pi)^{r_2}R_K}{w_K\sqrt{|d_K|}}.
\]
\end{theorem}
\begin{proof}
Choose an integral ideal $J$ in the class $C^{-1}$ (multiply any ideal of that class by a suitable integer). For $I\in C$
integral, $IJ$ is principal and contained in $J$: $IJ=\gamma\mathcal O$ with $\gamma\in J\setminus\{0\}$, determined up to a unit.
Conversely each $\gamma\in J\setminus\{0\}$ gives the integral ideal $I=\gamma J^{-1}\in C$. So the ideals to count correspond to the
$\mathcal O_K^\times$-orbits in $J\setminus\{0\}$, with $N(I)=|N(\gamma)|/N(J)$. Apply Theorem~\ref{thm:conecount} with $U=\mathcal
O_K^\times$, $\mathfrak a=J$, $\beta=0$, all signs, and $X=xN(J)$:
\[
\frac{2^{r_1}\pi^{r_2}R_K}{w_K\,2^{-r_2}\sqrt{|d_K|}N(J)}\,xN(J)+O\bigl(x^{1-1/n}\bigr)=\kappa_Kx+O\bigl(x^{1-1/n}\bigr).
\]
\end{proof}

The constant does not depend on the class: every class gets the same share of ideals. Figure~\ref{fig:idealcount} shows the two
classes of $\Z[\sqrt{-5}]$ racing each other along the line $\kappa x$, with $\kappa=\frac{2\pi}{2\sqrt{20}}=\frac\pi{\sqrt{20}}$.

\begin{figure}[ht]
\centering
\fig{ch08-idealcount}
\caption{The number of ideals of $\Z[\sqrt{-5}]$ of norm at most $x$, in the principal class (norms $a^2+5b^2$) and in the other class
(norms $2a^2+2ab+3b^2$), against $\kappa x$. At $x=20\,000$ the ratios to $x$ are $0.7023$ and $0.7024$, and $\kappa=0.70248\ldots$}
\label{fig:idealcount}
\end{figure}

For Chebotarev's theorem we need finer classes, defined by congruences. Fix an integer $m\ge1$. Let $I_m$ be the group of fractional
ideals in whose factorisation no prime dividing $m$ occurs. For $\alpha\in K^\times$ write $\alpha\equiv1\ (\operatorname{mod}^\times m)$ if $\alpha=a/b$
with $a,b\in\mathcal O_K$, $a\mathcal O$ and $b\mathcal O$ prime to $m$, and $a\equiv b\pmod{m\mathcal O}$; and $\alpha\gg0$ if
$\sigma(\alpha)>0$ for every real embedding $\sigma$. Let
\[
P_m^+=\{\alpha\mathcal O:\ \alpha\equiv1\ (\operatorname{mod}^\times m),\ \alpha\gg0\},\qquad \mathrm{Cl}_m^+=I_m/P_m^+ ,
\]
the \emph{narrow ray class group modulo $m$}. (Products and inverses of such $\alpha$ are again such, so $P_m^+$ is a subgroup.) For
$m=1$ it is the \emph{narrow class group}; it maps onto the class group, with kernel of order at most $2^{r_1}$.

\begin{lemma}\label{lem:mult}
If $\alpha\in\mathcal O_K$ and $\alpha\equiv1\ (\operatorname{mod}^\times m)$, then $\alpha\equiv1\pmod{m\mathcal O}$.
\end{lemma}
\begin{proof}
Write $\alpha=a/b$ as in the definition. Then $b(\alpha-1)=a-b\in m\mathcal O$. As $b\mathcal O+m\mathcal O=\mathcal O$, write
$1=xb+ym$; then $\alpha-1=x\,b(\alpha-1)+ym(\alpha-1)\in m\mathcal O$.
\end{proof}

\begin{theorem}[ideals in a ray class]\label{thm:raycount}
Every class $C\in\mathrm{Cl}_m^+$ contains integral ideals, and
\[
\#\{I\in C:\ I\subseteq\mathcal O_K,\ N(I)\le x\}=\kappa_{K,m}\,x+O\bigl(x^{1-1/n}\bigr),
\]
where $\kappa_{K,m}>0$ does not depend on $C$. In particular $\mathrm{Cl}_m^+$ is finite.
\end{theorem}
\begin{proof}
\emph{Integral representatives.} Let $I_0\in C$, and write $I_0=\mathfrak b\mathfrak c^{-1}$ with $\mathfrak b,\mathfrak c$ integral and
prime to $m$. The integer $N(\mathfrak c)$ lies in $\mathfrak c$ and is prime to $m$ (its prime factors lie under the primes of
$\mathfrak c$). So $a=N(\mathfrak c)^{\varphi(m)}$ is a positive integer, $a\equiv1\pmod m$, and $aI_0\in C$ is integral.

\emph{Counting.} Let $J$ be an integral ideal in $C^{-1}$, $U=\{u\in\mathcal O^\times:u\equiv1\pmod m,\ u\gg0\}$, and $w_U$, $R_U$ as
above. $U$ has finite index in $\mathcal O^\times$: it contains $\varepsilon^{2M}$ for every unit $\varepsilon$, where $M$ is the order
of $(\mathcal O/m\mathcal O)^\times$, and $\mathcal O^\times/(\mathcal O^\times)^{2M}$ is finite since $\mathcal O^\times$ is finitely
generated. For $I\in C$ integral, $IJ\in P_m^+$ is integral: $IJ=\gamma\mathcal O$ with $\gamma\in J$, $\gamma\gg0$,
$\gamma\equiv1\ (\operatorname{mod}^\times m)$, hence $\gamma\equiv1\pmod{m\mathcal O}$ by Lemma~\ref{lem:mult}; and two such generators differ by
an element of $U$. Since $J+m\mathcal O=\mathcal O$, the Chinese remainder theorem gives $\beta\in J$ with $\beta\equiv1\pmod
{m\mathcal O}$, and then $J\cap(1+m\mathcal O)=\beta+(J\cap m\mathcal O)=\beta+mJ$. So the ideals to count correspond to the
$U$-orbits of the totally positive elements of $\beta+mJ$, and Theorem~\ref{thm:conecount} (one sign pattern, $X=xN(J)$) gives
\[
\frac{\pi^{r_2}R_U}{w_U\,2^{-r_2}\sqrt{|d_K|}\,N(mJ)}\,xN(J)=\frac{(2\pi)^{r_2}R_U}{w_U\sqrt{|d_K|}\,m^n}\,x=:\kappa_{K,m}x
\]
up to $O(x^{1-1/n})$, since $N(mJ)=m^nN(J)$. The constant does not involve $C$.

\emph{Finiteness.} Distinct classes contain distinct ideals, and by Theorem~\ref{thm:idealcount} there are $O(x)$ integral ideals of
norm at most $x$ in all. If there were $h$ ray classes, then $h\kappa_{K,m}x\le O(x)$, so $h$ is bounded.
\end{proof}

\section{The analytic class number formula}

\begin{lemma}\label{lem:partialsum}
Let $(a_k)_{k\ge1}$ be complex numbers with $\sum_{k\le x}a_k=cx+O(x^\theta)$, where $0\le\theta<1$. Then $\sum a_kk^{-s}$ converges for
$\sigma>1$, and $\sum_ka_kk^{-s}-c\,\zeta(s)$ extends to a holomorphic function on $\sigma>\theta$.
\end{lemma}
\begin{proof}
Let $b_k=a_k-c$, so $B(x)=\sum_{k\le x}b_k=cx-c\lfloor x\rfloor+O(x^\theta)=O(x^\theta)$. By Abel summation, $\sum_{k\le y}b_kk^{-s}=
B(y)y^{-s}+s\int_1^yB(t)t^{-s-1}\dd t$. For $\sigma>\theta$ the first term tends to $0$ and the integral converges absolutely,
dominated on each compact set of the half-plane, where $|s|\le c'$ and $\sigma\ge\sigma_0>\theta$, by $c'C\,t^{\theta-\sigma_0-1}$; so $\sum b_kk^{-s}=s\int_1^\infty B(t)t^{-s-1}\dd t$ is
holomorphic there (Volume~4, Theorem~9.15(b)). For $\sigma>1$ we may add $c\,\zeta(s)$ term by term.
\end{proof}

\begin{theorem}[Dirichlet, Dedekind: the analytic class number formula]\label{thm:cnf}
$\zeta_K(s)$ extends to a meromorphic function on $\sigma>1-\frac1n$ whose only pole is a simple pole at $s=1$, with residue
\[
\operatorname*{Res}_{s=1}\zeta_K(s)=\frac{2^{r_1}(2\pi)^{r_2}\,h_K\,R_K}{w_K\sqrt{|d_K|}}.
\]
\end{theorem}
\begin{proof}
Let $a_k$ be the number of ideals of norm $k$. Summing Theorem~\ref{thm:idealcount} over the $h_K$ classes, $\sum_{k\le x}a_k=h_K\kappa_Kx+
O(x^{1-1/n})$. By Lemma~\ref{lem:partialsum}, $\zeta_K(s)-h_K\kappa_K\zeta(s)$ is holomorphic on $\sigma>1-\frac1n$, and $\zeta(s)-\frac1{s-1}$
is holomorphic on $\sigma>0$.
\end{proof}

Every invariant of Chapters~1 to~4 sits in one formula: signature, class number, regulator, roots of unity, discriminant. The class number
and the regulator appear only through the product $h_KR_K$, which is why they can trade size, as in real quadratic fields
(Chapter~4).

\begin{think}
For an imaginary quadratic field, $R_K=1$ and $w_K=2$ (unless $K=\Q(i)$ or $\Q(\sqrt{-3})$), so the residue is $\frac{\pi h_K}{\sqrt{|d_K|}}$.
Siegel proved in 1935 that $\log h_K\sim\frac12\log|d_K|$. What does this suggest about the size of the residue, and so about how
$\zeta_K$ behaves near $s=1$ as $|d_K|$ grows? (Volume~13 returns to this.)
\end{think}

\section{Abelian fields and Dirichlet \texorpdfstring{$L$}{L}-functions}

To use the formula one must compute the residue in another way. For subfields of cyclotomic fields, $\zeta_K$ factors into Dirichlet
$L$-functions, whose values at $1$ Dirichlet knew. We do it for subfields of $\Q(\zeta_p)$, which covers the cases we need.

Let $p$ be an odd prime, $G=\operatorname{Gal}(\Q(\zeta_p)/\Q)\cong(\Z/p\Z)^\times$ ($a\leftrightarrow\sigma_a$), $K\subseteq\Q(\zeta_p)$ a
subfield, $H=\operatorname{Gal}(\Q(\zeta_p)/K)$, and $X_K$ the set of Dirichlet characters $\chi$ modulo $p$ with $\chi=1$ on $H$. Then
$X_K$ is the dual of $G/H\cong\operatorname{Gal}(K/\Q)$, so $|X_K|=[K:\Q]$. Every $\chi\ne\chi_0$ modulo $p$ is primitive.

\begin{theorem}\label{thm:zetaab}
For $\sigma>1$,
\[
\zeta_K(s)=\zeta(s)\prod_{\chi\in X_K,\ \chi\ne\chi_0}L(s,\chi).
\]
Consequently $\displaystyle\prod_{\chi\in X_K,\ \chi\ne\chi_0}L(1,\chi)=\frac{2^{r_1}(2\pi)^{r_2}h_KR_K}{w_K\sqrt{|d_K|}}$.
\end{theorem}
\begin{proof}
Compare Euler factors, $n=[K:\Q]$. For $q\ne p$: $q$ is unramified in $\Q(\zeta_p)$, with Frobenius $\sigma_q$ (Theorem~\ref{thm:cyclosplit}),
hence unramified in $K$. If $P$ is a prime of $\Q(\zeta_p)$ above $q$, then $\sigma_q$ maps $\mathcal O_K$ to itself (as $K/\Q$ is Galois,
$G$ being abelian), fixes $P\cap K$, and acts as $x\mapsto x^q$ modulo it; so the Frobenius of $q$ in $K$ is the image $\bar q$ of $q$
in $G/H$. Let $f$ be its order: $q$ splits in $K$ into $g=n/f$ primes of norm $q^f$ (Theorems~\ref{thm:transitive} and~\ref{thm:decomp}),
and the Euler factor of $\zeta_K$ is $(1-q^{-fs})^{-g}$. On the other side, $\chi\mapsto\chi(q)$ is a homomorphism from $X_K$ onto
the group $\mu_f$ (a character of the cyclic group $\langle\bar q\rangle$ extends to $G/H$, because the restriction map
$\widehat{G/H}\to\widehat{\langle\bar q\rangle}$ has kernel the dual of $(G/H)/\langle\bar q\rangle$, of order $n/f$, hence image of order
$f$), with fibres of size $n/f$. As $\prod_{\xi\in\mu_f}(1-\xi T)=1-T^f$,
\[
\prod_{\chi\in X_K}\bigl(1-\chi(q)q^{-s}\bigr)^{-1}=\bigl(1-q^{-fs}\bigr)^{-n/f}.
\]
For $q=p$: $p$ is totally ramified in $\Q(\zeta_p)$ (Proposition~\ref{prop:lambda}), so $p-1=e(\Q(\zeta_p)/K)\,e(K/\Q)$ with each factor at
most the degree, and $p$ is totally ramified in $K$: one prime of norm $p$, Euler factor $(1-p^{-s})^{-1}$, which is the factor of
$\zeta(s)$, while $\chi(p)=0$ for the other characters. The Euler products converge absolutely, so they multiply. The last statement
is Theorem~\ref{thm:cnf}, since $\zeta(s)(s-1)\to1$ and the $L(s,\chi)$, $\chi\ne\chi_0$, are continuous at $1$ (Volume~5,
Theorem~3.6(b)).
\end{proof}

The theorem gives a new proof that $L(1,\chi)\ne0$ for every non-principal $\chi$ modulo $p$ (take $K=\Q(\zeta_p)$: the right side is
not zero), as promised in Chapter~5 for the even characters.

\begin{example}[the opening]\label{exa:h23}
$K=\Q(\sqrt{-23})$ is the quadratic subfield of $\Q(\zeta_{23})$ (Proposition~\ref{prop:gausssum}, $p^*=-23$), so $X_K=\{\chi_0,\chi\}$ with
$\chi=\bigl(\frac\cdot{23}\bigr)$, odd. Here $r_2=1$, $w_K=2$, $R_K=1$, $d_K=-23$, and Theorem~\ref{thm:zetaab} reads $L(1,\chi)=\frac{\pi h_K}
{\sqrt{23}}$. On the other hand, $\tau(\chi)=i\sqrt{23}$ (Gauss, Volume~4, Theorem~8.14) and $\sum_{a=1}^{22}\bigl(\frac a{23}\bigr)a=-69$,
so Dirichlet's formula gives $L(1,\chi)=-\frac{i\pi}{23\cdot i\sqrt{23}}(-69)=\frac{3\pi}{\sqrt{23}}$. Hence $h_K=3$. No coincidence:
the series of the opening \emph{is} $L(1,\chi)$, and it knows the class group.
\end{example}

\section{The class number of \texorpdfstring{$\Q(\zeta_p)$}{Q(zeta p)} and Kummer's criterion}

Let $p$ be an odd prime, $K=\Q(\zeta_p)$, $\zeta=\zeta_p$, $K^+=\Q(\eta)$ with $\eta=\zeta+\zeta^{-1}$ its maximal real subfield, and $n=p-1$.
We compare the class number formulas of $K$ and $K^+$. First the three invariants that differ.

\begin{proposition}\label{prop:kplus}
\begin{enumerate}[label=(\alph*)]
\item $\mathcal O_{K^+}=\Z[\eta]$, and $\mathcal O_K=\mathcal O_{K^+}\oplus\zeta\mathcal O_{K^+}$.
\item $|d_K|=p\,d_{K^+}^2$, hence $d_{K^+}=p^{(p-3)/2}$.
\item $w_K=2p$, $w_{K^+}=2$, and $R_K=2^{(p-3)/2}R_{K^+}$.
\end{enumerate}
\end{proposition}
\begin{proof}
(a) $\mathcal O_K=\Z[\zeta]$ (Theorem~\ref{thm:Zzeta}). As $\zeta^2=\eta\zeta-1$, the group $\Z[\eta]+\zeta\Z[\eta]$ is stable under multiplication
by $\zeta$, so it is a ring containing $\zeta$, hence equal to $\Z[\zeta]$; the sum is direct because $1,\zeta$ is a basis of $K$ over $K^+$.
If $x\in\mathcal O_{K^+}$, write $x=a+\zeta b$ with $a,b\in\Z[\eta]$; as $x\in K^+$, $b=0$. So $\mathcal O_{K^+}=\Z[\eta]$.

(b) Let $\omega_1,\dots,\omega_m$ ($m=n/2$) be an integral basis of $K^+$; by (a), the $\omega_i$ and $\zeta\omega_i$ form one of $K$. The
embeddings of $K$ come in pairs $\sigma,\bar\sigma$ restricting to the same embedding $\tau$ of $K^+$. In the matrix $(\sigma(\text{basis}))$,
the rows of $\sigma$ and $\bar\sigma$ are $\bigl(\tau(\omega_i),\ \sigma(\zeta)\tau(\omega_i)\bigr)$ and $\bigl(\tau(\omega_i),\
\sigma(\zeta)^{-1}\tau(\omega_i)\bigr)$. Subtracting the first from the second and reordering rows makes the matrix block triangular, with
diagonal blocks $T=(\tau(\omega_i))$ and $\Delta T$, $\Delta=\operatorname{diag}\bigl(\sigma(\zeta^{-1}-\zeta)\bigr)$. So $|d_K|=d_{K^+}^2
\prod_\sigma|\sigma(\zeta^{-1}-\zeta)|^2$, the product over one $\sigma$ in each pair, which is $|N_{K/\Q}(\zeta^{-1}-\zeta)|=|N(1-\zeta^2)|=p$.
With $|d_K|=p^{p-2}$ (Theorem~\ref{thm:Zzeta}) and $d_{K^+}>0$ ($K^+$ is totally real, Proposition~\ref{prop:sign}), $d_{K^+}=p^{(p-3)/2}$.

(c) $\mu_K=\{\pm\zeta^j\}$ (Proposition~\ref{prop:muKn}), and $K^+\subseteq\R$. By Kummer's lemma (Theorem~\ref{thm:kummerlemma}),
$\mathcal O_K^\times=\mu_K\,\mathcal O_{K^+}^\times$, so fundamental units $\varepsilon_1,\dots,\varepsilon_{m-1}$ of $K^+$ are fundamental units of
$K$. Each complex place of $K$ restricts to a real place of $K^+$, and $\ell_K(\varepsilon)=2\ell_{K^+}(\varepsilon)$ coordinate by
coordinate ($e_i=2$ for $K$, $1$ for $K^+$). An $(m-1)\times(m-1)$ minor is multiplied by $2^{m-1}$.
\end{proof}

For a character $\chi$ modulo $p$ define the \emph{generalised Bernoulli number} $B_{1,\chi}=\frac1p\sum_{a=1}^{p-1}\chi(a)a$. Let
$h^-=h_K/h_{K^+}$, the \emph{relative class number}.

\begin{theorem}[Kummer]\label{thm:hminus}
$\displaystyle h^-=2p\prod_{\chi\text{ odd}}\Bigl|\frac12B_{1,\chi}\Bigr|$, the product over the $\frac{p-1}2$ odd characters modulo $p$. Moreover
$\prod_{\chi\text{ odd}}B_{1,\chi}$ is a rational number, so $h^-=\pm2p\prod_{\chi\text{ odd}}\frac12B_{1,\chi}$.
\end{theorem}
\begin{proof}
$K^+$ is the fixed field of $\sigma_{-1}$, so $X_{K^+}$ is the set of even characters, and $X_K$ is the set of all characters. Dividing the
two formulas of Theorem~\ref{thm:zetaab} ($r_2=m$ for $K$, $r_1=m$ for $K^+$):
\[
\prod_{\chi\text{ odd}}L(1,\chi)=\frac{(2\pi)^mh_KR_K}{w_K\sqrt{|d_K|}}\cdot\frac{2\sqrt{d_{K^+}}}{2^mh_{K^+}R_{K^+}}
=h^-\,\frac{(2\pi)^m\,2^{m-1}\cdot2}{2p\cdot2^m\,p^{m/2}}=h^-\frac{\pi^m\,2^m}{2p\,p^{m/2}},
\]
by Proposition~\ref{prop:kplus}, since $\sqrt{|d_K|/d_{K^+}}=\sqrt{p\,d_{K^+}}=p^{m/2}$. By Dirichlet's formula and $|\tau(\overline\chi)|=\sqrt p$,
$|L(1,\chi)|=\frac\pi{p\sqrt p}\bigl|\sum_a\overline\chi(a)a\bigr|=\frac\pi{\sqrt p}|B_{1,\overline\chi}|$ for odd $\chi$, and $\chi\mapsto\overline\chi$
permutes the odd characters. Taking absolute values, $\pi^mp^{-m/2}\prod|B_{1,\chi}|=h^-\pi^m2^mp^{-m/2}/(2p)$, which is the formula.

The values of the characters lie in $\Q(\zeta_{p-1})$. An automorphism $\gamma$ of $\Q(\zeta_{p-1})$ maps $B_{1,\chi}$ to $B_{1,\gamma\circ\chi}$,
and $\gamma\circ\chi$ is again an odd character (it is a homomorphism, and $\gamma(-1)=-1$); so $\gamma$ permutes the factors of the product,
which is therefore fixed by the Galois group of $\Q(\zeta_{p-1})/\Q$, that is rational.
\end{proof}

To read off divisibility by $p$, we reduce the characters modulo a prime above $p$.

\begin{lemma}[the Teichm\"uller character]\label{lem:teich}
Let $\mathfrak p$ be a prime of $\Z[\zeta_{p-1}]$ above $p$. Then $\Z[\zeta_{p-1}]/\mathfrak p=\mathbb F_p$, and there is a unique character
$\omega$ modulo $p$ with values in $\mu_{p-1}$ such that $\omega(a)\equiv a\pmod{\mathfrak p}$ for all $a$. It generates the group of characters
modulo $p$, and $\omega(-1)=-1$.
\end{lemma}
\begin{proof}
$p\nmid p-1$ and $p\equiv1\pmod{p-1}$ has order $1$, so $p$ splits completely in $\Q(\zeta_{p-1})$ (Corollary~\ref{cor:cycloram}), and the
residue field is $\mathbb F_p$. The $p-1$ roots of $X^{p-1}-1$ in $\mu_{p-1}$ have distinct reductions modulo $\mathfrak p$, because
$\prod_{\xi\ne1}(1-\xi)=p-1\notin\mathfrak p$; so reduction is an isomorphism of groups $\mu_{p-1}\to\mathbb F_p^\times$, and $\omega$ is its
inverse composed with $a\mapsto a\bmod p$. It is injective, so it has order $p-1$ and generates the cyclic group of characters. Finally
$\omega(-1)^2=1$ and $\omega(-1)\equiv-1\not\equiv1$, so $\omega(-1)=-1$.
\end{proof}

So the odd characters are the $\omega^{j}$ with $j$ odd, $1\le j\le p-2$; that is, $\omega^{-1}$ (for $j=p-2$) and the $\omega^{k-1}$ with $k$
even, $2\le k\le p-3$. We write $x\equiv y\pmod p$ for $p$-integral rationals when $v_p(x-y)\ge1$, and similarly modulo $p^2$ and modulo
$\mathfrak p$.

\begin{lemma}[Bernoulli numbers modulo $p$]\label{lem:bernp}
Let $p\ge5$.
\begin{enumerate}[label=(\alph*)]
\item $v_p(B_j)\ge-1$ for every $j\ge0$.
\item For even $k$ with $2\le k\le p-3$, $B_k$ is $p$-integral and $\sum_{a=1}^{p-1}a^k\equiv pB_k\pmod{p^2}$.
\item (Voronoi) For such $k$ and every integer $c$ prime to $p$, with $q_a=\lfloor ca/p\rfloor$,
\[
(c^k-1)B_k\equiv k\,c^{k-1}\sum_{a=1}^{p-1}a^{k-1}q_a\pmod p.
\]
\end{enumerate}
\end{lemma}
\begin{proof}
Faulhaber's formula with $N=p$, using $\binom{j+1}i\frac1{j+1}=\binom ji\frac1{j+1-i}$, reads
\[
\sum_{a=0}^{p-1}a^j=pB_j+\sum_{l=2}^{j+1}\binom j{j+1-l}B_{j+1-l}\,\frac{p^l}l.
\]
(a) By induction on $j$, $B_0=1$. In the sum, $v_p(B_{j+1-l})\ge-1$ and $v_p(p^l/l)=l-v_p(l)\ge1$, so every term is $p$-integral, and so is
the left side; hence $v_p(pB_j)\ge0$.

(b) Take $j=k$. The term $l=2$ involves $B_{k-1}$, which is $0$ if $k\ge4$ ($k-1$ odd) and $-\frac12$ if $k=2$; in both cases the term has
$v_p\ge2$. For $l=3$, $v_p\ge-1+3=2$ as $p\ge5$; for $l\ge4$, $v_p\ge-1+l-v_p(l)\ge2$ since $p^{v_p(l)}\le l$ forces $v_p(l)\le l-3$. So
$\sum_{a=1}^{p-1}a^k\equiv pB_k\pmod{p^2}$. Let $c$ be a primitive root modulo $p$; $c^k\not\equiv1$ as $p-1\nmid k$, and $a\mapsto ca$
permutes the non-zero residues, so $(c^k-1)\sum a^k\equiv0\pmod p$. Hence $v_p(pB_k)\ge1$.

(c) Write $ca=pq_a+r_a$ with $1\le r_a\le p-1$; $a\mapsto r_a$ is a permutation. By the binomial theorem, $r_a^k=(ca-pq_a)^k\equiv
c^ka^k-k\,c^{k-1}a^{k-1}pq_a\pmod{p^2}$. Summing over $a$ and using (b) twice,
\[
pB_k\equiv c^kpB_k-k\,p\,c^{k-1}\sum_aa^{k-1}q_a\pmod{p^2},
\]
and we divide by $p$.
\end{proof}

\begin{theorem}[Kummer's criterion]\label{thm:kummercrit}
Let $p\ge5$ be prime. Then $p$ divides $h^-$ (that is, $v_p(h^-)\ge1$) if and only if $p$ divides the numerator of one of $B_2,B_4,\dots,B_{p-3}$.
In that case $p\mid h_p$: $p$ is irregular.
\end{theorem}
\begin{proof}
Let $\mathfrak p$ and $\omega$ be as in Lemma~\ref{lem:teich}. On $\Q$, the valuation $v_{\mathfrak p}$ is $v_p$, since $p$ is unramified in
$\Q(\zeta_{p-1})$. By Theorem~\ref{thm:hminus},
\[
h^-=\pm\frac{2}{2^{(p-1)/2}}\cdot\bigl(pB_{1,\omega^{-1}}\bigr)\cdot\prod_{\substack{2\le k\le p-3\\ k\text{ even}}}B_{1,\omega^{k-1}}.
\]
\emph{The factor $pB_{1,\omega^{-1}}$ is a $\mathfrak p$-unit}: $pB_{1,\omega^{-1}}=\sum_a\omega(a)^{-1}a\equiv\sum_aa^{-1}a=p-1\equiv-1\pmod{\mathfrak p}$.

\emph{The other factors.} Fix $k$, put $\chi=\omega^{k-1}$, and let $c$ be a primitive root modulo $p$. With $ca=pq_a+r_a$ as in
Lemma~\ref{lem:bernp}(c),
\[
pB_{1,\chi}=\sum_a\chi(r_a)r_a=\sum_a\chi(c)\chi(a)(ca-pq_a)=c\chi(c)\,pB_{1,\chi}-p\,\chi(c)\sum_a\chi(a)q_a,
\]
so $\bigl(1-c\chi(c)\bigr)B_{1,\chi}=-\chi(c)\sum_a\chi(a)q_a$, an exact identity. Now $c\chi(c)\equiv c^k\not\equiv1\pmod{\mathfrak p}$, so
$B_{1,\chi}$ is $\mathfrak p$-integral, and reducing modulo $\mathfrak p$ with $\chi(a)\equiv a^{k-1}$ and Lemma~\ref{lem:bernp}(c):
\[
(1-c^k)B_{1,\chi}\equiv-c^{k-1}\sum_aa^{k-1}q_a\equiv-(c^k-1)\frac{B_k}k\pmod{\mathfrak p},\qquad\text{so}\qquad B_{1,\omega^{k-1}}\equiv\frac{B_k}k\pmod{\mathfrak p}.
\]
Hence $v_p(h^-)=v_{\mathfrak p}(h^-)=\sum_kv_{\mathfrak p}(B_{1,\omega^{k-1}})$, a sum of non-negative terms, which is positive exactly when some
$B_k/k$ is $\equiv0\pmod{\mathfrak p}$, that is, as $B_k/k$ is a $p$-integral rational and $k<p$, when $p$ divides the numerator of $B_k$.
Finally $h_p=h^-h_{K^+}$ with $h_{K^+}\in\Z$, so $v_p(h_p)\ge v_p(h^-)$.
\end{proof}

For $p=37$: $B_{32}=-\frac{7\,709\,321\,041\,217}{510}$ and $7\,709\,321\,041\,217=37\cdot683\cdot305\,065\,927$. So $37\mid h_{37}$, closing the
signature exercise of Chapter~5 (in fact $h_{37}=h^-=37$ and $h_{37}^+=1$). The table of Theorem~\ref{thm:hminus} up to $41$:
\[
h^-=1\ (p\le19),\quad 3\ (23),\quad 8\ (29),\quad 9\ (31),\quad 37\ (37),\quad 121\ (41),
\]
and a direct computation verifies every value against PARI's class group computations.

\begin{remark}
$h^-$ is always an integer, because $h_{K^+}$ divides $h_K$ (a consequence of class field theory, Chapter~10). The converse half of Kummer's
criterion, \emph{if $p\mid h_p$ then $p\mid h^-$}, is Theorem~5.34 in Washington's book; with it, Theorem~\ref{thm:kummercrit} becomes the full
criterion of Chapter~5. Herbrand (1932) and Ribet (1976) refined it: $p\mid B_k$ if and only if a specific eigenspace of the $p$-part of the
class group, cut out by $\omega^{1-k}$, is non-zero. Ribet's proof built the missing piece of the class group out of a modular form, a first
glimpse of Chapter~14.
\end{remark}

\begin{think}
Kummer's tables of 1850 stopped at $p=163$. Which irregular primes did he find? Before computing them, guess from the heuristic that
each of the $\frac{p-3}2$ numbers $B_k$ is divisible by $p$ ``with probability $\frac1p$''. What proportion of primes should be regular?
\end{think}

\section{Densities and Frobenius elements}

From now on we work with \emph{relative} extensions $L/K$ of number fields. For a prime $\mathfrak p$ of $K$, write $\mathfrak p\mathcal O_L=
\prod\mathfrak P_i^{e_i}$, $k_{\mathfrak p}=\mathcal O_K/\mathfrak p$, $f_i=[\mathcal O_L/\mathfrak P_i:k_{\mathfrak p}]$, so that $N\mathfrak P_i=
N\mathfrak p^{f_i}$. The results of Chapter~2 extend.

\begin{proposition}\label{prop:relative}
Let $L/K$ be an extension of number fields of degree $d$.
\begin{enumerate}[label=(\alph*)]
\item $\sum_ie_if_i=d$.
\item Only finitely many $\mathfrak p$ ramify in $L$.
\end{enumerate}
Suppose now $L/K$ Galois with group $G$, and let $\mathfrak P$ be a prime of $L$ above $\mathfrak p$.
\begin{enumerate}[label=(\alph*),resume]
\item $G$ acts transitively on the primes above $\mathfrak p$, which have common $e$ and $f$, with $efg=d$.
\item The decomposition group $D_{\mathfrak P}$ maps onto $\operatorname{Gal}(k_{\mathfrak P}/k_{\mathfrak p})$, the inertia group $I_{\mathfrak P}$ has
order $e$, and if $\mathfrak p$ is unramified there is a unique $\operatorname{Frob}_{\mathfrak P}\in D_{\mathfrak P}$ with
$\operatorname{Frob}_{\mathfrak P}(x)\equiv x^{N\mathfrak p}\pmod{\mathfrak P}$; it generates $D_{\mathfrak P}$, which has order $f$. Moreover
$\operatorname{Frob}_{g\mathfrak P}=g\operatorname{Frob}_{\mathfrak P}g^{-1}$, so $\mathfrak p$ has a Frobenius class
$\operatorname{Frob}_{\mathfrak p}$, and $\mathfrak p$ splits completely if and only if $\operatorname{Frob}_{\mathfrak p}=\{1\}$.
\item If $K\subseteq E\subseteq L$ and $\mathfrak q=\mathfrak P\cap E$, then $\operatorname{Frob}_{\mathfrak P}(L/E)=\operatorname{Frob}_{\mathfrak
P}(L/K)^{f(\mathfrak q/\mathfrak p)}$. If moreover $E/K$ is Galois, $\operatorname{Frob}_{\mathfrak P}(L/K)$ restricts to
$\operatorname{Frob}_{\mathfrak q}(E/K)$.
\end{enumerate}
\end{proposition}
\begin{proof}
(a) For $a\in K^\times$, $N_{L/\Q}(a)=N_{K/\Q}(a)^d$ (each embedding of $K$ extends in $d$ ways to $L$), so $N(a\mathcal O_L)=N(a\mathcal O_K)^d$.
Let $h=h_K$; then $\mathfrak p^h=a\mathcal O_K$, and by multiplicativity of the norm
\[
N\mathfrak p^{\,hd}=N(a\mathcal O_L)=N(\mathfrak p\mathcal O_L)^h=\prod_iN\mathfrak P_i^{\,e_ih}=N\mathfrak p^{\,h\sum e_if_i}.
\]
(b) If $\mathfrak P$ is ramified over $\mathfrak p$, it is ramified over $p=\mathfrak p\cap\Z$ in $L/\Q$ (ramification indices multiply in
the tower $\Q\subseteq K\subseteq L$), so $p\mid d_L$ (Theorem~\ref{thm:ramdisc}).

(c), (d) The proofs of Theorems~\ref{thm:transitive} and~\ref{thm:decomp} apply word for word, with $\Z$, $p$, $\mathbb F_p$ replaced by
$\mathcal O_K$, $\mathfrak p$, $k_{\mathfrak p}$: in the first, $N_{L/K}(x)\in\mathfrak Q\cap\mathcal O_K=\mathfrak p\subseteq\mathfrak P$; in
the second, the polynomial $\prod_{\sigma\in G}(X-\sigma(a))$ has coefficients in $\mathcal O_K$, and $\operatorname{Gal}(k_{\mathfrak
P}/k_{\mathfrak p})$ is cyclic, generated by $y\mapsto y^{N\mathfrak p}$ (Volume~2, Chapter~9). The conjugation formula is proved as in
Chapter~2.

(e) $\operatorname{Frob}_{\mathfrak P}(L/E)$ is the element of $D_{\mathfrak P}(L/E)\subseteq D_{\mathfrak P}(L/K)$ acting as $x\mapsto
x^{N\mathfrak q}=x^{N\mathfrak p^f}$ on $k_{\mathfrak P}$, and so is $\operatorname{Frob}_{\mathfrak P}(L/K)^f$; uniqueness in (d) gives equality.
If $E/K$ is Galois, $\operatorname{Frob}_{\mathfrak P}(L/K)$ maps $\mathcal O_E$ to itself, fixes $\mathfrak q$ and acts as $x\mapsto
x^{N\mathfrak p}$ modulo $\mathfrak q$.
\end{proof}

\begin{definition}
A set $S$ of primes of $K$ has \emph{Dirichlet density} $\delta$ if
\[
\lim_{s\to1^+}\frac{\sum_{\mathfrak p\in S}N\mathfrak p^{-s}}{\log\frac1{s-1}}=\delta ,
\]
the limit over real $s>1$. The \emph{lower} and \emph{upper} densities are the $\liminf$ and $\limsup$.
\end{definition}

\begin{proposition}\label{prop:primesum}
As $s\to1^+$, $\sum_{\mathfrak p}N\mathfrak p^{-s}=\log\frac1{s-1}+O(1)$, and the primes $\mathfrak p$ with $N\mathfrak p$ not a prime number
contribute $O(1)$. So the set of all primes has density $1$, and changing a set by finitely many primes, or by primes of residue degree at
least $2$ over $\Q$, does not change its upper or lower density.
\end{proposition}
\begin{proof}
By Proposition~\ref{prop:zetaK}, $\log\zeta_K(s)=\sum_{\mathfrak p}N\mathfrak p^{-s}+\sum_{\mathfrak p}\sum_{k\ge2}\frac{N\mathfrak p^{-ks}}k$ for
real $s>1$. A prime $p$ has at most $n$ primes above it, so the double sum and the sum over the $\mathfrak p$ with $N\mathfrak p=p^f$, $f\ge2$,
are both at most $n\sum_p\sum_{k\ge2}p^{-k}=n\sum_p\frac1{p(p-1)}<\infty$. And $\log\zeta_K(s)=\log\frac1{s-1}+O(1)$ by
Theorem~\ref{thm:cnf}.
\end{proof}

\begin{theorem}[Kronecker, 1880]\label{thm:kronecker}
If $L/K$ is Galois of degree $d$, the primes of $K$ that split completely in $L$ have density $\frac1d$. In particular, if all but finitely many
primes of $K$ split completely in $L$, then $L=K$.
\end{theorem}
\begin{proof}
Apply Proposition~\ref{prop:primesum} to $L$. Up to $O(1)$ we may discard the primes $\mathfrak P$ of $L$ above ramified primes, and those with
$f(\mathfrak P/\mathfrak p)\ge2$ (their norm $N\mathfrak p^f$ is not prime). A remaining $\mathfrak P$ has $e=f=1$ over $\mathfrak p$, so all
primes above $\mathfrak p$ do (Proposition~\ref{prop:relative}(c)): $\mathfrak p$ splits completely, into $d$ primes of norm $N\mathfrak p$.
Hence $\log\frac1{s-1}+O(1)=d\sum_{\mathfrak p\text{ split}}N\mathfrak p^{-s}$.
\end{proof}

Figure~\ref{fig:cubefrob} tests the general statement on $X^3-2$, whose splitting field $L=\Q(\sqrt[3]2,\zeta_3)$ has group $\mathfrak S_3$. By
Dedekind's theorem (Theorem~\ref{thm:dedekindcycle}), the number of roots of $X^3-2$ modulo $p$ is the number of fixed points of
$\operatorname{Frob}_p$ acting on the three cube roots.

\begin{think}
Guess, before reading the caption, the proportion of primes $p$ for which $X^3-2$ has exactly one root modulo $p$. (Hint: for $p\equiv2\pmod3$,
cubing is a bijection of $\mathbb F_p$.) Then guess the proportion for which it has three.
\end{think}

\begin{figure}[ht]
\centering
\fig{ch08-cubefrob}
\caption{For the $17\,982$ primes $3<p<2\cdot10^5$: the proportion of $p$ for which $X^3-2$ has three, one or no roots modulo $p$ (dark bars),
against the proportion of identity, transpositions and $3$-cycles in $\mathfrak S_3$ (light bars): $\frac16$, $\frac12$, $\frac13$.}
\label{fig:cubefrob}
\end{figure}

\section{Primes with prescribed norms modulo \texorpdfstring{$m$}{m}}

Let $F$ be a number field of degree $n_F$, $m\ge1$, $\zeta=e^{2\pi i/m}$ and $M=F(\zeta)$. An automorphism $\sigma\in\operatorname{Gal}(M/F)$ is
determined by $\sigma(\zeta)=\zeta^a$, and $\sigma\mapsto a\bmod m$ is an injective homomorphism into $(\Z/m\Z)^\times$; let $A$ be its image.
This is Dirichlet's setting with $\Q$ replaced by $F$.

\begin{lemma}\label{lem:cyclofrob}
Let $\mathfrak p$ be a prime of $F$ not dividing $m$. Then $\mathfrak p$ is unramified in $M$, and $\operatorname{Frob}_{\mathfrak p}$
corresponds to $N\mathfrak p\bmod m$. In particular $N(I)\bmod m\in A$ for every $I\in I_m$.
\end{lemma}
\begin{proof}
Let $\mathfrak P\mid\mathfrak p$ and $\sigma\in D_{\mathfrak P}$ with $\sigma(x)\equiv x^{N\mathfrak p}\pmod{\mathfrak P}$
(Proposition~\ref{prop:relative}(d)). If $\sigma(\zeta)=\zeta^a$, then $\zeta^a\equiv\zeta^{N\mathfrak p}\pmod{\mathfrak P}$. The $m$-th roots
of unity are distinct modulo $\mathfrak P$, since $\prod_{\xi\ne1}(1-\xi)=m\notin\mathfrak P$; so $a\equiv N\mathfrak p\pmod m$. The same argument
shows that an element of the inertia group, which acts trivially on $\mathcal O_M/\mathfrak P$, fixes $\zeta$, hence is the identity: $e=1$.
\end{proof}

\begin{theorem}[Dirichlet's theorem over a number field]\label{thm:dirichletF}
For every $a\in A$, the primes $\mathfrak p\nmid m$ of $F$ with $N\mathfrak p\equiv a\pmod m$ have density $\frac1{|A|}$. Equivalently, each element of
$\operatorname{Gal}(F(\zeta_m)/F)$ is the Frobenius of a set of primes of density $\frac1{[F(\zeta_m):F]}$.
\end{theorem}
\begin{proof}
\emph{Step 1: norms fill $A$.} Let $B\subseteq A$ be the subgroup generated by the $N\mathfrak p\bmod m$, $\mathfrak p\nmid m$, and $E\subseteq M$
the fixed field of the corresponding subgroup of $\operatorname{Gal}(M/F)$. Every $\mathfrak p\nmid m$ has Frobenius in $B$, so its Frobenius in
$E/F$, the restriction (Proposition~\ref{prop:relative}(e)), is trivial: $\mathfrak p$ splits completely in $E$. By Kronecker's theorem, $E=F$, so
$B=A$.

\emph{Step 2: $L$-functions.} For a character $\psi$ of $A$ and $I\in I_m$ put $\psi_N(I)=\psi(N(I)\bmod m)$. It is trivial on $P_m^+$: if
$\alpha=a/b\equiv1\ (\operatorname{mod}^\times m)$ and $\alpha\gg0$, then $N(\alpha\mathcal O)=N_{F/\Q}(\alpha)=N_{F/\Q}(a)/N_{F/\Q}(b)$, and $N_{F/\Q}(a)\equiv N_{F/\Q}(b)
\pmod m$ because $\sigma(a)\equiv\sigma(b)\pmod{m\mathcal O_{\tilde F}}$ for every embedding $\sigma$ into the Galois closure $\tilde F$, and
$m\mathcal O_{\tilde F}\cap\Z=m\Z$. So $\psi_N$ is a character of the finite group $\mathrm{Cl}_m^+(F)$, non-trivial if $\psi\ne1$ by Step~1. Let
\[
L(s,\psi_N)=\sum_{I\subseteq\mathcal O_F,\ I\in I_m}\frac{\psi_N(I)}{N(I)^s}=\prod_{\mathfrak p\nmid m}\bigl(1-\psi(N\mathfrak p)N\mathfrak p^{-s}\bigr)^{-1}
\qquad(\sigma>1),
\]
the Euler product as in Proposition~\ref{prop:zetaK}. For $\psi\ne1$, group the ideals by ray class: $L(s,\psi_N)=\sum_C\psi_N(C)Z_C(s)$, where
by Theorem~\ref{thm:raycount} and Lemma~\ref{lem:partialsum} each $Z_C(s)-\kappa_{F,m}\zeta(s)$ is holomorphic on $\sigma>1-\frac1{n_F}$, and
$\sum_C\psi_N(C)=0$. So $L(s,\psi_N)$ is holomorphic on $\sigma>1-\frac1{n_F}$. For $\psi=1$, $L(s,1_N)=\zeta_F(s)\prod_{\mathfrak p\mid m}(1-
N\mathfrak p^{-s})$.

\emph{Step 3: factorisation.} A prime $\mathfrak p\nmid m$ splits in $M$ into $|A|/f$ primes of norm $N\mathfrak p^f$, where $f$ is the order of
$N\mathfrak p$ in $A$ (Lemma~\ref{lem:cyclofrob} and Proposition~\ref{prop:relative}(d)). As in the proof of Theorem~\ref{thm:zetaab},
$\prod_{\psi\in\widehat A}(1-\psi(N\mathfrak p)T)=(1-T^f)^{|A|/f}$. Hence, for $\sigma>1$,
\[
\zeta_M(s)=E(s)\prod_{\psi\in\widehat A}L(s,\psi_N),\qquad E(s)=\prod_{\mathfrak P\mid m}\bigl(1-N\mathfrak P^{-s}\bigr)^{-1},
\]
a finite product, holomorphic and non-zero on $\sigma>0$.

\emph{Step 4: non-vanishing.} By Theorem~\ref{thm:cnf} for $M$ and $F$, $(s-1)\zeta_M(s)$ and $(s-1)L(s,1_N)$ have non-zero limits at $s=1$. So
$\prod_{\psi\ne1}L(s,\psi_N)$ has a finite non-zero limit; each factor is holomorphic at $1$, so none vanishes there.

\emph{Step 5: densities.} For $\psi\ne1$ and real $s\in\left]1,2\right]$, $S_\psi(s)=\sum_{\mathfrak p\nmid m}\psi(N\mathfrak p)N\mathfrak p^{-s}$ plus the
absolutely convergent sum $T_\psi(s)$ of the terms $k\ge2$ is a continuous logarithm of $L(s,\psi_N)$ (Proposition~\ref{prop:zetaK}). Since
$L(s,\psi_N)\to L(1,\psi_N)\ne0$, near which a holomorphic logarithm exists, $S_\psi(s)$ stays bounded as $s\to1^+$. By orthogonality,
\[
\sum_{\substack{\mathfrak p\nmid m\\ N\mathfrak p\equiv a\ (m)}}N\mathfrak p^{-s}=\frac1{|A|}\sum_{\psi\in\widehat A}\overline{\psi(a)}S_\psi(s)
=\frac1{|A|}\log\frac1{s-1}+O(1).
\]
\end{proof}

For $F=\Q$ this is Dirichlet's theorem in Dirichlet density, with a new proof of $L(1,\chi)\ne0$ in which the real characters cause no trouble: the
simple pole of $\zeta_{\Q(\zeta_m)}$ does the work that the positivity tricks of Volume~5 did. The price is the counting of ideals in
$\Q(\zeta_m)$.

\section{Chebotarev's density theorem}

\begin{theorem}[Chebotarev, 1922]\label{thm:chebotarev}
Let $L/K$ be a Galois extension of number fields with group $G$, and $C\subseteq G$ a conjugacy class. The primes $\mathfrak p$ of $K$, unramified in
$L$, with $\operatorname{Frob}_{\mathfrak p}=C$ have density $\frac{|C|}{|G|}$.
\end{theorem}

The proof has two reductions. The first, due to Deuring, moves from a conjugacy class to an element of a cyclic group. For $\sigma\in G$ of
order $f$, let $E=L^{\langle\sigma\rangle}$, so that $L/E$ is cyclic with group $\langle\sigma\rangle$, and let $C$ be the class of $\sigma$ in
$G$.

\begin{lemma}[Deuring]\label{lem:deuring}
For real $s>1$,
\[
\sum_{\substack{\mathfrak q\text{ of }E\text{, unramified in }L\\ \operatorname{Frob}_{\mathfrak q}(L/E)=\sigma}}N\mathfrak q^{-s}
=\frac{|C_G(\sigma)|}f\sum_{\substack{\mathfrak p\text{ of }K\text{, unramified in }L\\ \operatorname{Frob}_{\mathfrak p}=C}}N\mathfrak p^{-s}+O(1),
\]
where $C_G(\sigma)$ is the centraliser. Consequently, if the left side has density $\frac1f$, the primes of $K$ with Frobenius class $C$ have
density $\frac1{|C_G(\sigma)|}=\frac{|C|}{|G|}$.
\end{lemma}
\begin{proof}
On the left, the primes $\mathfrak q$ of residue degree $f(\mathfrak q/\mathfrak p)\ge2$ over $K$ contribute $O(1)$ (their norm is not prime), as
do those above primes $\mathfrak p$ ramified in $L$. Let $\mathcal Q$ be the set of the remaining $\mathfrak q$: $f(\mathfrak q/\mathfrak
p)=1$, $\mathfrak p$ unramified, $\operatorname{Frob}_{\mathfrak q}(L/E)=\sigma$. Let $\mathcal P$ be the set of primes $\mathfrak P$ of $L$, above
unramified primes of $K$, with $\operatorname{Frob}_{\mathfrak P}(L/K)=\sigma$.

If $\mathfrak q\in\mathcal Q$ and $\mathfrak P\mid\mathfrak q$, then $\operatorname{Frob}_{\mathfrak P}(L/K)=\operatorname{Frob}_{\mathfrak
P}(L/E)=\sigma$ by Proposition~\ref{prop:relative}(e), as $f(\mathfrak q/\mathfrak p)=1$; and $\mathfrak P$ is the only prime above $\mathfrak q$,
since $D_{\mathfrak P}(L/E)=\langle\sigma\rangle$ is everything. Conversely, if $\mathfrak P\in\mathcal P$, then $D_{\mathfrak P}(L/K)=\langle
\sigma\rangle=\operatorname{Gal}(L/E)$, so $f(\mathfrak P/\mathfrak q)=|D_{\mathfrak P}(L/E)|=f=f(\mathfrak P/\mathfrak p)$, hence
$f(\mathfrak q/\mathfrak p)=1$ and $\operatorname{Frob}_{\mathfrak q}(L/E)=\sigma$. So $\mathfrak P\mapsto\mathfrak P\cap E$ is a bijection
$\mathcal P\to\mathcal Q$ preserving norms.

Now fix $\mathfrak p$ with $\operatorname{Frob}_{\mathfrak p}=C$, and $\mathfrak P_0\mid\mathfrak p$ with Frobenius $\sigma$. The primes above
$\mathfrak p$ are the $g\mathfrak P_0$, with $g\mathfrak P_0=g'\mathfrak P_0$ iff $g^{-1}g'\in D_{\mathfrak P_0}=\langle\sigma\rangle$, and
$\operatorname{Frob}_{g\mathfrak P_0}=g\sigma g^{-1}$ equals $\sigma$ iff $g\in C_G(\sigma)$. As $\langle\sigma\rangle\subseteq C_G(\sigma)$, exactly
$|C_G(\sigma)|/f$ primes of $\mathcal P$ lie above $\mathfrak p$, each of norm $N\mathfrak p$. Summing gives the formula. The last
statement is the orbit--stabiliser count $|C|=|G|/|C_G(\sigma)|$.
\end{proof}

So it suffices to prove the theorem for cyclic extensions. The second reduction is Chebotarev's: cross with a cyclotomic field, so that a cyclic
extension becomes cyclotomic over a suitable subfield, where Theorem~\ref{thm:dirichletF} applies.

\begin{lemma}\label{lem:disjoint}
Let $L/K$ be cyclic of degree $d$, and $q_1,\dots,q_t$ distinct primes with $q_i\equiv1\pmod d$ and $q_i\nmid d_Kd_L$; put $m=q_1\cdots q_t$ and
$M=K(\zeta_m)$. Then $\operatorname{Gal}(M/K)\cong(\Z/m\Z)^\times\cong\prod_i(\Z/q_i\Z)^\times$, and $L\cap M=K$, so that
$\operatorname{Gal}(LM/K)\cong\operatorname{Gal}(L/K)\times\operatorname{Gal}(M/K)$.
\end{lemma}
\begin{proof}
$k=K\cap\Q(\zeta_m)$ is unramified over $\Q$ at every $p\nmid m$ (being inside $\Q(\zeta_m)$, Corollary~\ref{cor:cycloram}) and at the $q_i$
(being inside $K$, and $q_i\nmid d_K$). By Minkowski (Theorem~\ref{thm:dkbig} with Theorem~\ref{thm:ramdisc}), $k=\Q$, and
Galois theory gives $\operatorname{Gal}(M/K)\cong\operatorname{Gal}(\Q(\zeta_m)/\Q)$.

Let $\mathfrak P$ be a prime of $M$ above $q_i$. In $\Q(\zeta_m)$, $q_i$ has ramification index $q_i-1$ (it is totally ramified in
$\Q(\zeta_{q_i})$ and unramified in $\Q(\zeta_{m/q_i})$), and $q_i$ is unramified in $K$, so $e(\mathfrak P/\mathfrak P\cap K)\ge q_i-1$. On
the other hand the inertia group $I_{\mathfrak P}\subseteq\operatorname{Gal}(M/K)$ acts trivially on $K(\zeta_{m/q_i})$, in which the primes above
$q_i$ are unramified (Lemma~\ref{lem:cyclofrob}); so $I_{\mathfrak P}$ lies in the factor $(\Z/q_i\Z)^\times$, of order $q_i-1$, and by
Proposition~\ref{prop:relative}(d) it is that whole factor. Now $E=L\cap M$ is Galois over $K$ and unramified at every prime ($L/K$ is
unramified above the $q_i$, $M/K$ elsewhere), so every $I_{\mathfrak P}$ acts trivially on $E$: an element of $I_{\mathfrak P}$ restricts to an
element of the inertia group of $\mathfrak P\cap E$ in $E/K$, which is trivial. The $I_{\mathfrak P}$ generate $\prod_i(\Z/q_i\Z)^\times$, so
$E=K$. The last statement is Galois theory for the compositum.
\end{proof}

\begin{proof}[Proof of Theorem~\ref{thm:chebotarev}]
By Lemma~\ref{lem:deuring} applied to $\sigma$ and $E=L^{\langle\sigma\rangle}$, it suffices to show: \emph{if $L/K$ is cyclic of degree $d$, every
$\tau\in G=\operatorname{Gal}(L/K)$ is the Frobenius of a set of primes of density $\frac1d$.} Write $P_\tau$ for that set, and $\underline\delta$,
$\overline\delta$ for lower and upper densities.

Take $M$ as in Lemma~\ref{lem:disjoint} (primes $q\equiv1\pmod d$ not dividing $d_Kd_L$ exist in any number, by Dirichlet's theorem), and
$A=\operatorname{Gal}(M/K)$. Call $\rho\in A$ \emph{good} if the order of $\tau$ divides the order of $\rho$. For good $\rho$, let $\Sigma=\langle
(\tau,\rho)\rangle\subseteq G\times A=\operatorname{Gal}(LM/K)$, of order $\operatorname{ord}\rho$, and $F=(LM)^\Sigma$. Then $\Sigma\cap
(G\times\{1\})=\{1\}$: if $\rho^k=1$ then $\operatorname{ord}\rho\mid k$, so $\tau^k=1$. Since $G\times\{1\}=\operatorname{Gal}(LM/M)$, the
group $\operatorname{Gal}(LM/FM)=\Sigma\cap\operatorname{Gal}(LM/M)$ is trivial: $LM=FM=F(\zeta_m)$. By Theorem~\ref{thm:dirichletF}, the primes of
$F$ with Frobenius $(\tau,\rho)$ in $LM/F$ have density $\frac1{|\Sigma|}$. As $\operatorname{Gal}(LM/K)$ is abelian, Lemma~\ref{lem:deuring}
(with $L,E$ replaced by $LM,F$, centraliser of order $d|A|$) shows that the primes $\mathfrak p$ of $K$ with $\operatorname{Frob}_{\mathfrak
p}(LM/K)=(\tau,\rho)$ have density $\frac{|\Sigma|}{d|A|}\cdot\frac1{|\Sigma|}=\frac1{d|A|}$. These sets are disjoint for different $\rho$ and
contained in $P_\tau$ (restriction of Frobenius, Proposition~\ref{prop:relative}(e)), so
\[
\underline\delta(P_\tau)\ge\frac1d\cdot\frac{\#\{\rho\in A\text{ good}\}}{|A|}.
\]
Now let $t\to\infty$. A $\rho=(\rho_1,\dots,\rho_t)\in\prod(\Z/q_i\Z)^\times$ is good as soon as some $\rho_i$ has order divisible by $d$. In the
cyclic group $(\Z/q_i\Z)^\times$, of order divisible by $d$, the elements whose order is divisible by $d$ form a proportion $c\ge\frac{\varphi(d)}d>0$
(an element $x$ of $\Z/N\Z$, $d\mid N$, has order divisible by $d$ iff for each prime $\ell\mid d$, $v_\ell(x)\le v_\ell(N)-v_\ell(d)$; the
proportion is $\prod_{\ell\mid d}(1-\ell^{-(v_\ell(N)-v_\ell(d)+1)})\ge\prod_{\ell\mid d}(1-\frac1\ell)$). So the proportion of bad $\rho$ is at
most $(1-c)^t\to0$, and $\underline\delta(P_\tau)\ge\frac1d$ for every $\tau\in G$. Finally the $P_\tau$ partition the unramified primes, of
density $1$, so $\overline\delta(P_\tau)\le1-\sum_{\tau'\ne\tau}\underline\delta(P_{\tau'})\le1-\frac{d-1}d=\frac1d$.
\end{proof}

\begin{remark}
Chebotarev's theorem also holds with \emph{natural} density, $\#\{N\mathfrak p\le x:\operatorname{Frob}_{\mathfrak p}=C\}\sim\frac{|C|}{|G|}\frac x
{\log x}$, by the same reductions plus a zero-free region for Hecke $L$-functions (the prime number theorem of Volume~5 in this setting). Lagarias
and Odlyzko (1977) made it effective: under GRH, the least prime with a given Frobenius is $O((\log|d_L|)^2)$. Artin used the crossing trick of
the proof above to prove his reciprocity law in 1927 (Chapter~10).
\end{remark}

\begin{corollary}\label{cor:chebcons}
\begin{enumerate}[label=(\alph*)]
\item Every element of a Galois group $G$ is the Frobenius of infinitely many primes.
\item If $f\in\Z[X]$ is monic irreducible with Galois group $G$ acting on its $n$ roots, the proportion of primes modulo which $f$ factors with
degrees $(d_1,\dots,d_r)$ is the proportion of elements of $G$ with that cycle type.
\item If $L/K$ and $L'/K$ are Galois and the same primes (up to a set of density $0$) split completely in both, then $L=L'$.
\end{enumerate}
\end{corollary}
\begin{proof}
(a) A set of positive density is infinite. (b) Theorem~\ref{thm:dedekindcycle} and Proposition~\ref{prop:frobembed}: for $p\nmid\operatorname{disc}(f)$
the factorisation pattern is the cycle type of $\operatorname{Frob}_p$. (c) is Exercise~\ref{ex:bauer}.
\end{proof}

Figure~\ref{fig:cycletypes} of Chapter~2 and Figure~\ref{fig:cubefrob} are instances of (b).

\section{Exercises}

\begin{exercise}[$\star$]\label{ex:hminusp}
Let $p\equiv3\pmod4$, $p>3$. Show that $h(\Q(\sqrt{-p}))=-\frac1p\sum_{a=1}^{p-1}\bigl(\frac ap\bigr)a$. Compute it for $p=7,23,47$, and deduce that
there are more quadratic residues than non-residues in $\left]0,\frac p2\right[$.
\end{exercise}

\begin{exercise}[$\star\star$]\label{ex:sqrt2}
Let $K=\Q(\sqrt2)$ and $\chi_8$ the character modulo $8$ with $\chi_8(\pm1)=1$, $\chi_8(\pm3)=-1$. Show that $\zeta_K(s)=\zeta(s)L(s,\chi_8)$, compute
$L(1,\chi_8)$ with Dirichlet's formula for even characters (Volume~5, Theorem~3.10), and deduce $h_K=1$ from the class number formula.
\end{exercise}

\begin{exercise}[$\star\star$]\label{ex:twosq}
Show that $\zeta_{\Q(i)}(s)=\zeta(s)L(s,\chi_{-4})$, and deduce that the number $r_2(n)$ of pairs $(a,b)\in\Z^2$ with $a^2+b^2=n$ is
$4\sum_{d\mid n}\chi_{-4}(d)$. Recover from Theorem~\ref{thm:idealcount} that $\sum_{n\le x}r_2(n)=\pi x+O(\sqrt x)$, Gauss's circle problem.
\end{exercise}

\begin{exercise}[$\star\star$]\label{ex:splitdensity}
Let $K$ be a number field, not necessarily Galois, and $L$ its Galois closure. Show that the primes $p$ that split completely in $K$ are exactly those
that split completely in $L$, so they have density $\frac1{[L:\Q]}$. For $K=\Q(\sqrt[3]2)$ this is $\frac16$, not $\frac13$.
\end{exercise}

\begin{exercise}[$\star\star$]\label{ex:rootevery}
\begin{enumerate}[label=(\alph*)]
\item Show that $(X^2-2)(X^2-3)(X^2-6)$ has a root modulo every prime, but no rational root.
\item Let $f\in\Z[X]$ be irreducible of degree $n\ge2$. Show that $f$ has no root modulo $p$ for a set of primes of positive density. (Jordan: a
transitive group on $n\ge2$ points contains an element without fixed points.)
\end{enumerate}
\end{exercise}

\begin{exercise}[$\star\star$]\label{ex:cycliccubic}
Let $f=X^3-3X+1$, with roots $2\cos\frac{2\pi k}9$. Show that its splitting field is the real cubic subfield of $\Q(\zeta_9)$, that $f$ has three roots
modulo $p\ne3$ if $p\equiv\pm1\pmod9$ and none otherwise, and check the densities $\frac13$, $\frac23$ against Chebotarev.
\end{exercise}

\begin{exercise}[$\star\star\star$]\label{ex:bauer}
\emph{Bauer's theorem.} Let $L,L'$ be Galois extensions of $K$. Show that if every prime of $K$ outside a set of density $0$ that splits completely in
$L$ also splits completely in $L'$, then $L'\subseteq L$. Deduce Corollary~\ref{cor:chebcons}(c): a Galois extension is determined by its split
primes. (So the quadratic field $\Q(\sqrt d)$ is determined by the primes $p$ with $\bigl(\frac dp\bigr)=1$, which quadratic reciprocity describes by
congruences; Chapter~10 asks which fields are described by congruences.)
\end{exercise}

\begin{beauty}
Counting lattice points in a cone is the most naive of computations, yet it produces a formula in which every invariant of a number field
appears at once, and Dirichlet's infinite series turn out to know the class group. Kummer squeezed from it a link between ideals and Bernoulli
numbers that is still a research subject (Iwasawa theory, the Herbrand--Ribet theorem). And Chebotarev's trick of crossing with a cyclotomic field,
invented to count primes, became in Artin's hands the proof of the reciprocity law that governs all abelian extensions: the density theorem is
the bridge between the analytic world of this chapter and the algebraic world of class field theory.
\end{beauty}

\section{Hints}

Try each exercise for a while before reading its hint, and read the hint before the solution.

\begin{hint}{ex:hminusp}
$\Q(\sqrt{-p})$ is the quadratic subfield of $\Q(\zeta_p)$; follow Example~\ref{exa:h23}. For the last part, split the sum at $\frac p2$ as in
Volume~5, Exercise~3.6.
\end{hint}
\begin{hint}{ex:sqrt2}
The splitting of odd primes in $\Q(\sqrt2)$ depends on $\bigl(\frac2p\bigr)$; and $2=(\sqrt2)^2$ ramifies. The fundamental unit is $1+\sqrt2$, and
$\cot\frac\pi8=1+\sqrt2$.
\end{hint}
\begin{hint}{ex:twosq}
Compare Euler factors at $2$, at $p\equiv1$ and at $p\equiv3\pmod4$, then count elements instead of ideals ($w=4$).
\end{hint}
\begin{hint}{ex:splitdensity}
A prime splits completely in $K$ iff all its Frobenius elements in $L$ fix every embedding of $K$ (Proposition~\ref{prop:frobembed}).
\end{hint}
\begin{hint}{ex:rootevery}
(a) The product of two non-residues is a residue. (b) Count elements with a fixed point with Burnside's lemma: the average number of fixed
points of a transitive group is $1$.
\end{hint}
\begin{hint}{ex:cycliccubic}
$2\cos\frac{2\pi}9=\zeta_9+\zeta_9^{-1}$; the Galois group of the real subfield is $(\Z/9\Z)^\times/\{\pm1\}$.
\end{hint}
\begin{hint}{ex:bauer}
Compare the densities of primes splitting completely in $L$ and in the compositum $LL'$.
\end{hint}

\section{Solutions}

\begin{solution}{ex:hminusp}
\idea Example~\ref{exa:h23} works for every $p\equiv3\pmod4$: the class number formula on one side, Dirichlet's formula on the other.
\writeup Since $p^*=-p$, $K=\Q(\sqrt{-p})$ is the quadratic subfield of $\Q(\zeta_p)$ (Proposition~\ref{prop:gausssum}), and Theorem~\ref{thm:zetaab}
gives $L(1,\chi)=\frac{2\pi h}{2\sqrt p}=\frac{\pi h}{\sqrt p}$ for the Legendre symbol $\chi$, which is odd; here $w=2$ as $p>3$, $R=1$ and $|d_K|=p$
($-p\equiv1\pmod4$). By Gauss, $\tau(\chi)=i\sqrt p$, and Dirichlet's formula gives $L(1,\chi)=-\frac{i\pi}{p\cdot i\sqrt p}\sum_a\chi(a)a=
-\frac\pi{p^{3/2}}\sum_a\chi(a)a$. Comparing, $h=-\frac1p\sum_a\chi(a)a$. For $p=7$: $\sum=1+2-3+4-5-6=-7$, $h=1$. For $p=23$: $h=3$
(Example~\ref{exa:h23}). For $p=47$ the sum is $-235$ and $h=5$. Finally let $S=\sum_{a=1}^{p-1}\chi(a)a=-ph$, $S_1=\sum_{a<p/2}\chi(a)a$ and
$T=\sum_{a<p/2}\chi(a)$, the number of residues minus the number of non-residues in $\left]0,\frac p2\right[$. Pairing $a$ with $p-a$ and using
$\chi(-1)=-1$ gives $S=2S_1-pT$. Splitting $S$ into even terms $b=2a$ and odd terms $b=p-2a$ ($a<\frac p2$) gives $S=2\chi(2)S_1-\chi(2)(pT-2S_1)
=\chi(2)(4S_1-pT)$. Eliminating $S_1$: $S=\chi(2)(2S+pT)$, so $T=\frac{(\chi(2)-2)S}p=\bigl(2-\bigl(\frac2p\bigr)\bigr)h>0$.
\end{solution}

\begin{solution}{ex:sqrt2}
\idea Same as Theorem~\ref{thm:zetaab}, with $8$ in place of $p$: the splitting law of $\Q(\sqrt2)$ is a congruence modulo $8$.
\writeup For an odd prime $p$, $p$ splits in $\Q(\sqrt2)$ iff $2$ is a square modulo $p$ (Theorem~\ref{thm:kd} with $X^2-2$), iff $p\equiv\pm1
\pmod8$, iff $\chi_8(p)=1$; otherwise $p$ is inert. So the Euler factor of $\zeta_K$ at $p$ is $(1-p^{-s})^{-2}$ or $(1-p^{-2s})^{-1}$, which is
$(1-p^{-s})^{-1}(1-\chi_8(p)p^{-s})^{-1}$ in both cases. At $2$, $2\mathcal O=(\sqrt2)^2$, factor $(1-2^{-s})^{-1}$, and $\chi_8(2)=0$. So
$\zeta_K=\zeta\,L(\cdot,\chi_8)$. Now $\chi_8$ is even and primitive, with $\tau(\chi_8)=\zeta_8+\zeta_8^7-\zeta_8^3-\zeta_8^5=2\cos\frac\pi4-2\cos
\frac{3\pi}4=2\sqrt2$, so
\[
L(1,\chi_8)=-\frac1{2\sqrt2}\Bigl(2\log\sin\frac\pi8-2\log\sin\frac{3\pi}8\Bigr)=\frac1{\sqrt2}\log\cot\frac\pi8=\frac{\log(1+\sqrt2)}{\sqrt2}.
\]
The class number formula, with $r_1=2$, $w=2$, $R=\log(1+\sqrt2)$ and $d_K=8$, gives $L(1,\chi_8)=\frac{4hR}{2\sqrt8}=\frac{h\log(1+\sqrt2)}{\sqrt2}$.
So $h=1$.
\end{solution}

\begin{solution}{ex:twosq}
\idea Euler factors again, then the count of ideals becomes a count of lattice points because $\Z[i]$ is principal.
\writeup $p\equiv1\pmod4$ splits, $p\equiv3$ is inert, $2=-i(1+i)^2$ ramifies (Volume~1); as in Exercise~\ref{ex:sqrt2}, $\zeta_{\Q(i)}=\zeta\,L(\cdot,
\chi_{-4})$. Comparing coefficients of $n^{-s}$, the number of ideals of norm $n$ is $\sum_{d\mid n}\chi_{-4}(d)$. Every ideal is principal, generated
by $w=4$ elements $a+bi$ of norm $a^2+b^2$, so $r_2(n)=4\sum_{d\mid n}\chi_{-4}(d)$. Theorem~\ref{thm:idealcount} with $h=1$, $r_2=1$, $R=1$, $w=4$,
$|d|=4$ gives $\kappa=\frac{2\pi}{4\cdot2}=\frac\pi4$, so $\sum_{n\le x}r_2(n)=4(\frac\pi4x+O(\sqrt x))=\pi x+O(\sqrt x)$: the lattice points in the
disc of radius $\sqrt x$. (The truth is believed to be $O(x^{1/4+\varepsilon})$, an open problem.)
\end{solution}

\begin{solution}{ex:splitdensity}
\idea Splitting completely is a statement about all embeddings at once, and the embeddings of $K$ generate $L$.
\writeup Let $p$ be unramified in $L$ (all but finitely many). By Proposition~\ref{prop:frobembed}, $p$ splits completely in $K$ iff the Frobenius
$\sigma$ of some (any) prime of $L$ above $p$ fixes every embedding $\tau:K\to L$, that is, $\sigma\tau=\tau$ for all $\tau$. Then $\sigma$ fixes every
$\tau(K)$, and these fields generate $L$, so $\sigma=1$: $p$ splits completely in $L$. The converse is clear. By Kronecker's theorem the density is
$\frac1{[L:\Q]}$. For $K=\Q(\sqrt[3]2)$, $[L:\Q]=6$: only the $p\equiv1\pmod3$ for which $2$ is a cube, one prime in six; compare the first bar of
Figure~\ref{fig:cubefrob}.
\end{solution}

\begin{solution}{ex:rootevery}
\idea (a) The three numbers $2,3,6$ cannot all be non-residues. (b) Count permutations without fixed points.
\writeup (a) There are no rational roots. Modulo $p=2$ or $3$, $X^2-2$ or $X^2-3$ has the root $0$. For $p\ge5$, $\bigl(\frac6p\bigr)=\bigl(\frac2p\bigr)
\bigl(\frac3p\bigr)$, so if $2$ and $3$ are non-residues, $6$ is a residue. (b) Let $G$ act transitively on the $n$ roots. Burnside's lemma (Volume~2, Theorem~3.6):
the average number of fixed points is the number of orbits, $1$. The identity has $n\ge2$ fixed points, so some element has none; let $C$ be its class.
By Chebotarev and Dedekind (Corollary~\ref{cor:chebcons}(b)), for a set of primes of density $\frac{|C|}{|G|}>0$, $\operatorname{Frob}_p\in C$ has no
fixed root, so $f$ has no linear factor modulo $p$. In (a), the polynomial is reducible, which is how it escapes.
\end{solution}

\begin{solution}{ex:cycliccubic}
\idea The roots are $\zeta+\zeta^{-1}$, $\zeta=\zeta_9^k$, so everything happens inside $\Q(\zeta_9)$.
\writeup With $c=\zeta_9+\zeta_9^{-1}=2\cos\frac{2\pi}9$: $c^3=\zeta_9^3+\zeta_9^{-3}+3c=-1+3c$, since $\zeta_9^3+\zeta_9^{-3}=2\cos\frac{2\pi}3=-1$. So
$c^3-3c+1=0$, and the other roots are $\sigma_2(c)$, $\sigma_4(c)$. So the splitting field is $\Q(c)=\Q(\zeta_9)^+$, cyclic of degree $3$, with group
$(\Z/9\Z)^\times/\{\pm1\}$. For $p\ne3$ (the discriminant is $81$), the Frobenius is the class of $p$; it is trivial iff $p\equiv\pm1\pmod9$, and then $f$
has three roots; otherwise it is a $3$-cycle and $f$ has none. Chebotarev predicts $\frac13$ and $\frac23$, which is also Dirichlet's theorem for the
two classes $\pm1$ among the six units modulo $9$.
\end{solution}

\begin{solution}{ex:bauer}
\idea If $L'\not\subseteq L$, the compositum is bigger than $L$, so fewer primes split completely in it.
\writeup Let $N=LL'$, Galois over $K$. A prime unramified in $N$ splits completely in $N$ iff it splits completely in $L$ and in $L'$ (its Frobenius in
$N$ is trivial iff its restrictions to $L$ and $L'$ are, as $\operatorname{Gal}(N/K)$ embeds into $\operatorname{Gal}(L/K)\times
\operatorname{Gal}(L'/K)$). By hypothesis, up to density $0$ the primes splitting completely in $L$ split completely in $N$. Kronecker's theorem gives
$\frac1{[N:K]}=\frac1{[L:K]}$, so $N=L$ and $L'\subseteq L$. For (c) apply this both ways.
\end{solution}

\section*{Sources and further reading}

\begin{itemize}
\item P.~G.~L.~Dirichlet, ``Recherches sur diverses applications de l'analyse infinit\'esimale \`a la th\'eorie des nombres'', \emph{J.\ reine angew.\
Math.}\ 19 (1839) and 21 (1840). The class number formula for quadratic forms.
\item E.~E.~Kummer, ``Allgemeiner Beweis des Fermat'schen Satzes \dots'', \emph{J.\ reine angew.\ Math.}\ 40 (1850), 130--138. The criterion and the
first irregular primes.
\item S.~Lang, \emph{Algebraic Number Theory}, 2nd ed., Springer GTM~110, 1994, Chapters~VI and~VIII. Counting ideals in cones, the Lipschitz
principle, and Chebotarev.
\item L.~C.~Washington, \emph{Introduction to Cyclotomic Fields}, 2nd ed., Springer GTM~83, 1997, Chapters~4 and~5. Generalised Bernoulli numbers,
the $h^-$ formula, Kummer's criterion in full.
\item P.~Stevenhagen and H.~W.~Lenstra, Jr., ``Chebotar\"ev and his density theorem'', \emph{Math.\ Intelligencer} 18 (1996), 26--37. The history, and
the proof followed here.
\item N.~Tschebotareff, ``Die Bestimmung der Dichtigkeit einer Menge von Primzahlen, welche zu einer gegebenen Substitutionsklasse geh\"oren'',
\emph{Math.\ Ann.}\ 95 (1926), 191--228.
\item J.~C.~Lagarias and A.~M.~Odlyzko, ``Effective versions of the Chebotarev density theorem'', in \emph{Algebraic Number Fields}, Academic Press,
1977, 409--464.
\item K.~A.~Ribet, ``A modular construction of unramified $p$-extensions of $\Q(\mu_p)$'', \emph{Invent.\ Math.}\ 34 (1976), 151--162.
\item \emph{History.} Frobenius's conjecture of 1880 and its publication in 1896, and Chebotarev's life in Odessa, are
in the Stevenhagen--Lenstra article above.
\end{itemize}
